\section*{Chapitre \ref{ch:b3:conservation-biology} --- Biologie de la conservation}

\begin{solution}{exo:b3:conservation-biology:1}
\omterm{def:b3:conservation-biology:small}{Stochasticité démographique}~: le hasard dans les naissances, les morts
et les sexes de quelques individus --- le kakapo à cinquante et un, avec
plus de mâles que de femelles. \omterm{def:b3:conservation-biology:small}{Stochasticité environnementale}~: une
mauvaise année qui frappe tous les individus à la fois --- une
sécheresse, une maladie, un hiver rude qui s'abat sur une population de
quelques centaines. \omterm{def:b3:conservation-biology:small}{Effet Allee}~: la croissance par tête qui baisse à
faible densité --- la tourte voyageuse, dont les colonies ne pouvaient
plus se reproduire une fois raréfiées, et dont les prédateurs restants
n'étaient plus submergés.
\end{solution}

\begin{solution}{exo:b3:conservation-biology:2}
La taille d'une population idéale --- constante, à panmixie, à sexes
équilibrés, à tailles de familles poissonniennes --- qui perdrait de
l'hétérozygotie au même rythme que la population réelle. Des nombres
inégaux de mâles et de femelles reproducteurs~; des fluctuations
d'effectif, qui pèsent du côté des générations les plus petites~; une
variance des tailles de familles supérieure à celle de Poisson, avec
quelques individus qui produisent l'essentiel des jeunes (ainsi que la
non-panmixie et le chevauchement des générations).
\end{solution}

\begin{solution}{exo:b3:conservation-biology:3}
$S = cA^{z}$. Fraction conservée $= (A'/A)^{z}$~: $0.5^{0.25} = 0.84$~;
$0.1^{0.25} = 0.56$~; $0.01^{0.25} = 0.32$ --- un centième de l'aire
garde encore un tiers des espèces, à terme.
\end{solution}

\begin{solution}{exo:b3:conservation-biology:4}
$N_{e} \ge 50$ tient la montée de la consanguinité au-dessous de $1/100$
par génération, le rythme auquel les éleveurs jugent la dépression
tolérable~: cela protège de la perte de valeur sélective à court terme.
$N_{e} \ge
500$ équilibre la perte de variation par dérive et son apport par
mutation, de sorte que la population garde la variation quantitative
dont elle a besoin pour s'adapter~: cela protège le long terme. Les
populations réelles ont besoin de cinq à dix fois ces nombres en
effectif recensé.
\end{solution}

\begin{solution}{exo:b3:conservation-biology:5}
$N_{e} = 4\times 8\times 120/128 = 30$~; l'hétérozygotie baisse de $1/60 =
\qty{1.7}{\%}$ par génération~; le quart est perdu quand $(1 - 1/60)^{t} =
0.75$, $t = \ln 0.75/\ln(59/60) = 17$ générations.
\end{solution}

\begin{solution}{exo:b3:conservation-biology:6}
$1/N_{e} = \tfrac{1}{5}(4/1000 + 1/20) = 0.0108$, $N_{e} = 93$~: une
génération à vingt a fait se comporter cinq générations comme
quatre-vingt-treize. Hétérozygotie conservée $(1 - 1/186)^{5} = 0.973$,
contre $(1 -
1/2000)^{5} = 0.9975$ pour un millier constant.
\end{solution}

\begin{solution}{exo:b3:conservation-biology:7}
$F = 1 - (1 - 1/20)^{5} = 0.23$~; valeur sélective en baisse de
$5\times 2.3 = \qty{11}{\%}$. Tout descendant d'une femelle du Texas et
d'un mâle de Floride a $F = 0$~; si les huit nouvelles venues font les
deux cinquièmes des femelles reproductrices, les deux cinquièmes de la
génération suivante sont hors consanguinité et le $F$ moyen tombe à
environ $0.6\times 0.23 = 0.14$ en une génération --- ce qu'a montré le
rétablissement.
\end{solution}

\begin{solution}{exo:b3:conservation-biology:8}
$r = \ln(100/31)/8 = \qty{0.15}{yr^{-1}}$. Poursuivi jusqu'en 2020~:
$100\,\mathrm{e}^{0.15\times 17} \approx 1200$. Le parc en tenait une
centaine~: les territoires se sont remplis, les wapitis ont décliné, les
meutes se sont entretuées, et la gale et la maladie de Carré se sont
répandues --- la dépendance à la densité s'est installée une fois le
territoire vide occupé.
\end{solution}

\begin{solution}{exo:b3:conservation-biology:9}
(a) Un déclin de \qty{60}{\%} en dix ans dépasse \qty{50}{\%}~: en
danger. (b) $180 < 250$ individus matures~: en danger critique. (c) Aire
de répartition \qty{3000}{km^2}, au-dessous de \qty{5000}{}~: en danger.
(d) \num{30000} individus et un déclin de \qty{10}{\%}~: aucun seuil
atteint --- non menacé (au plus quasi menacé, à surveiller).
\end{solution}

\begin{solution}{exo:b3:conservation-biology:10}
Les $N$ couples échouent tous avec la probabilité $(1/2)^{N}$~: $0.25$,
$0.031$, $0.001$, $10^{-6}$. Le risque démographique baisse
exponentiellement avec $N$ et devient négligeable au-delà de quelques
dizaines d'individus~; au-dessus, le risque qui reste est
environnemental, qui ne baisse pas si vite, parce qu'il frappe tous les
individus ensemble.
\end{solution}

\begin{solution}{exo:b3:conservation-biology:11}
Un fragment d'un dixième~: $0.1^{0.25} = 0.56$ des espèces initiales.
Dix fragments en abritent entre $0.56$ (s'ils abritent tous les mêmes
espèces) et, en principe, plus que le tout (si chacun en abritait de
différentes --- impossible au-delà du pool régional). Là où l'on tombe
dépend du renouvellement entre fragments~: des habitats différents dans
des fragments différents poussent le total vers le haut~; les espèces
qui ont besoin de grandes aires, d'habitat intérieur, ou qui sont
perdues dans chaque fragment par les mêmes \omterm{prop:b3:conservation-biology:area}{effets de lisière} le poussent
vers le bas~; et la dette d'extinction de chaque fragment se paie
séparément, de sorte que le total à long terme est d'ordinaire
au-dessous de celui du tout.
\end{solution}

\begin{solution}{exo:b3:conservation-biology:12}
La dérive retire de l'hétérozygotie au rythme $1/2N_{e}$ par
génération~; les immigrants remplacent une fraction $m$ du pool génique
par des allèles venus d'ailleurs. Avec $mN_{e} \approx 1$,
$m \approx 1/N_{e}$, soit le double du rythme de perte, de sorte que la
variation est renouvelée plus vite qu'elle ne décroît et que la
population reste proche de la source en fréquences alléliques
($F_{ST} = 1/(1
+ 4N_{e}m) \approx 0.2$). Le coût~: les allèles localement favorisés
sont dilués de $m$ par génération, de sorte que l'adaptation locale ne
survit que pour des traits dont le coefficient de sélection dépasse
environ $1/N_{e}$~; les différences locales faiblement sélectionnées
sont submergées.
\end{solution}

\begin{solution}{pb:b3:conservation-biology:1}
\textbf{1.} $N_{e} = 4\times 20\times 40/60 = 53$, contre un recensement de 60.
\textbf{2.} $1/N_{e} = \tfrac{1}{5}(1/500 + 1/100 + 1/60 + 1/51 + 1/60) =
0.0130$, $N_{e} = 77$~; les générations d'étranglement à 51 et 60
dominent, celle de 500 ne compte presque pas.
\textbf{3.} $F = 1 - (1 - 1/154)^{5} = 0.032$~; en linéaire $5/154 = 0.032$.
\textbf{4.} $6\times 0.32 = \qty{1.9}{\%}$.
\textbf{5.} $(1 - 1/154)^{5} = 0.968$~: \qty{96.8}{\%} conservés.
\textbf{6.} Dix générations à $N_{e} = 53$~: incrément $1 - (1 -
1/107)^{10} = 0.090$, total $F = 1 - 0.968\times 0.910 = 0.12$~; perte de
valeur sélective $6\times 1.2 = \qty{7}{\%}$~; un siècle.
\textbf{7.} Les couples local--local font $(60/70)^{2} = 0.73$ des
accouplements et leurs descendants gardent $F \approx 0.03$ (plus un
petit incrément)~; tous les autres ont $F = 0$~: en moyenne
$F \approx 0.73\times 0.035 = 0.026$, une baisse d'un quart en une
génération.
\textbf{8.} $F$ est l'identité par descendance au sein d'un individu, et
tout descendant hors consanguinité le remet à zéro~; l'hétérozygotie est
une propriété des fréquences alléliques, et les allèles des migrants ne
se répandent qu'à mesure qu'ils sont transmis. Dans une population
encore petite, la dérive les retirera de nouveau en quelques dizaines de
générations, de sorte que le sauvetage doit être répété ou la population
agrandie.
\textbf{9.} $N_{e} = N$ pour des sexes équilibrés~: 500 adultes. Depuis
70 à $r =
0.05$~: $\ln(500/70)/0.05 = 39$ ans.
\textbf{10.} La population idéale a des tailles de familles
poissonniennes, de variance égale à la moyenne de deux~;
$N_{e} = 4N/(2 + V_{k})$, de sorte que faire contribuer chaque couple
également ($V_{k} = 0$) donne $N_{e} = 2N$~: aucune lignée n'est perdue
par la chance du nid.
\textbf{11.} Recueillir et conserver son sperme pour une insémination
artificielle, et l'employer sur plusieurs femelles non apparentées au
fil des ans --- le coût est que ses allèles pourraient en venir à
dominer un petit pool, de sorte qu'il faut plafonner sa part près de
$1/N_{e}$~; ou cryoconserver du tissu pour un clonage ou une production
de gamètes ultérieurs, au prix du délai et du risque technique. Ne rien
faire perd pour de bon les allèles de la lignée.
\textbf{12.} \qty{160}{} millions de dollars pour quelque 430 oiseaux~:
environ \num{370000} par oiseau. Cher par perroquet, bon marché pour une
espèce --- et moins que quelques kilomètres d'autoroute, qui est la
comparaison que font réellement les budgets.
\textbf{13.} $(300/2000)^{0.25} = 0.62$~: environ 112 espèces conservées
sur 180, une dette d'extinction de 68.
\textbf{14.} Chaque fragment $(100/2000)^{0.25}\times 180 = 85$ espèces.
Entièrement différentes~: jusqu'à 255, impossible au-delà du pool de
180~; identiques~: 85. Ce qui décide, c'est à quel point les habitats
des fragments diffèrent et combien d'espèces peuvent persister dans
\qty{100}{km^2}~; en pratique entre 85 et 112.
\textbf{15.} $(320/2000)^{0.25}\times 180 = 114$, contre 85 pour trois
fragments isolés identiques~: le corridor laisse la recolonisation et le
flux de gènes faire agir trois petites populations comme une plus grande.
\textbf{16.} $N_{e} = 50$ avec des sexes équilibrés fait 25 couples~:
\qty{1250}{km^2}. Les \qty{300}{km^2} entiers abritent six couples,
$N_{e} \approx 12$~; un fragment, deux. Aucun ne suffit~: le prédateur
est perdu quel que soit le compte d'espèces, et protéger l'aire dont il
a besoin protège tout ce qui est plus petit --- une espèce parapluie.
\textbf{17.} $\ln(200/10)/0.03 = 100$ ans~: la dette se paie sur un
siècle, longtemps après que la forêt a été coupée.
\textbf{18.} Les lisières retirent de l'habitat intérieur à chaque
fragment, de sorte que l'aire utilisable est moindre que l'aire
cartographiée et la perte plus grande que ne le prédit la loi~; une
matrice de reprise forestière ou de haies dont certaines espèces peuvent
s'accommoder ajoute de l'aire efficace et rend la perte plus faible.
\textbf{19.} $r = \ln(100/31)/8 = \qty{0.146}{yr^{-1}}$~; temps de
doublement $\ln 2/0.146 = 4.7$ ans.
\textbf{20.} $N(8) = 120/[1 + (120/31 - 1)\mathrm{e}^{-1.17}] = 120/(1 +
2.87\times 0.31) = 63$. Les 100 observés sont la valeur exponentielle
elle-même --- $r$ a été ajusté sur ces deux points mêmes --- face à la
logistique 63~: dans les premières années les loups n'ont rencontré
aucune dépendance à la densité --- des territoires vides et des wapitis
abondants.
\textbf{21.} $\ln(4000/\num{17000})/20 = -0.072$~: \qty{7}{\%} par an.
Cent loups qui prennent \num{2000} wapitis par an sur une harde de
\num{8000} en moyenne font un quart par an --- plus qu'il n'en faut pour
mener le déclin à eux seuls, même si la chasse hors du parc, la
sécheresse et la prédation des grizzlys sur les faons y ont eu leur part.
\textbf{22.} Soixante individus matures, au-dessous de 250~: en danger
critique. Le prédateur à $N_{e} \approx 12$ (quelques dizaines
d'animaux)~: en danger critique également.
\textbf{23.} De $3\times 10^{9}$ à $1$ en 44 ans~: $r = -\ln(3\times
10^{9})/44 = -0.50$ par an, soit une division par deux tous les
dix-sept mois. Aucune récolte ne retire la moitié de milliards chaque
année~; la chute fut plus lente d'abord et bien plus rapide à la fin,
parce qu'un nicheur colonial au-dessous d'une densité seuil cesse de se
reproduire tandis que ses prédateurs ne sont plus submergés --- l'\omterm{def:b3:conservation-biology:small}{effet
Allee} a changé un déclin en effondrement.
\textbf{24.} En 1850 la population se comptait en milliards, son taux de
croissance mesuré était positif et sa variance petite, et une analyse de
viabilité projette la dynamique qu'on lui donne~: elle aurait mis le
risque à zéro. Elle aurait manqué le forçage extérieur --- les chemins
de fer et le télégraphe qui laissaient les chasseurs suivre et vider
chaque dortoir, le défrichement --- et le seuil d'Allee chez un nicheur
colonial, dont ni l'un ni l'autre n'est dans la démographie d'une espèce
abondante~; une AVP extrapole, et les causes d'extinction sont
d'ordinaire nouvelles.
\textbf{25.} $N_{e} \approx 53$ aujourd'hui~; $F \approx 0.12$ après dix
générations de plus à 60 (\qty{7}{\%} de la valeur sélective)~; la forêt
rétrécie garde environ \qty{62}{\%} de ses espèces.
\end{solution}
